% ============================================================
% TGP v1 — Preambuła
% Teoria Generowanej Przestrzeni: materia jako źródło przestrzeni
% ============================================================

\usepackage[utf8]{inputenc}
\usepackage[T1]{fontenc}
\usepackage{newunicodechar}
\newunicodechar{✓}{\ensuremath{\checkmark}}
\newunicodechar{✗}{\ensuremath{\times}}
% Fallbacks dla \cmark / \xmark (pifont-free)
\providecommand{\cmark}{\ensuremath{\checkmark}}
\providecommand{\xmark}{\ensuremath{\times}}
\newunicodechar{△}{\ensuremath{\triangle}}
\newunicodechar{★}{\ensuremath{\bigstar}}
\newunicodechar{∝}{\ensuremath{\propto}}
\newunicodechar{φ}{\ensuremath{\varphi}}
\newunicodechar{ψ}{\ensuremath{\psi}}
\newunicodechar{σ}{\ensuremath{\sigma}}
\newunicodechar{Φ}{\ensuremath{\varPhi}}
\newunicodechar{₀}{\ensuremath{_0}}
\newunicodechar{₁}{\ensuremath{_1}}
\newunicodechar{₂}{\ensuremath{_2}}
\newunicodechar{₃}{\ensuremath{_3}}
\newunicodechar{₄}{\ensuremath{_4}}
\usepackage{tikz}
\usetikzlibrary{arrows,arrows.meta,positioning,calc}
\usepackage[polish]{babel}
\usepackage{amsmath,amssymb,amsfonts,amsthm}
\usepackage{mathtools}   % \xleftrightarrow, \xrightarrow z etykietą
% Fallback dla \slashed{D} (operator Diraca) -- uzywany w sek09
\providecommand{\slashed}[1]{\not #1}
\usepackage[hidelinks]{hyperref}
% Fix: \k{} (polskie znaki) w tytulach sekcji nie psuje bookmarkow PDF
\pdfstringdefDisableCommands{\let\k\@gobble}
\usepackage{geometry}
% --- Extended appendix numbering (A, B, ..., Z, AA, AB, ...) ---
% TGP v1 has >26 appendices; default \Alph{section} overflows at A-Z.
% alphalph.sty provides \AlphAlph which extends naturally to AA, AB, ...
\usepackage{alphalph}
\makeatletter
\let\tgp@origappendix\appendix
\renewcommand{\appendix}{%
  \tgp@origappendix
  \renewcommand{\thesection}{\AlphAlph{\value{section}}}%
}
\makeatother
\usepackage{longtable,booktabs,tabularx,array}
\usepackage[shortlabels]{enumitem}
\usepackage{xcolor}
\geometry{margin=2.3cm}

% --- Status twierdzeń TGP ---
\newcommand{\statuslabel}[1]{\textbf{[#1]}}  % etykieta logiczna: [Aksjomat], [Hipoteza], ...

% --- Symbole TGP ---
\renewcommand{\Phi}{\varPhi}            % pole generujące przestrzeń
\newcommand{\PhiZero}{\varPhi_0}        % wartość próżniowa / odniesienia
\newcommand{\PhiEff}{\varPhi_{\mathrm{eff}}}   % efektywna wartość pola
\newcommand{\Nzero}{N_0}               % absolutna nicość
\newcommand{\Ffield}{\mathcal{F}}       % funkcja stanu (miara asymetrii)
\newcommand{\Floc}{\mathcal{F}_{\mathrm{loc}}}
\newcommand{\Srel}{S_{\mathrm{rel}}}    % entropia relacyjna
\newcommand{\GN}{G_{\mathrm{N}}}
\newcommand{\lP}{\ell_{\mathrm{P}}}
\newcommand{\Zp}{\mathbb{Z}_2}

% --- Otoczenia (angielskie nazwy) ---
\newtheorem{axiom}{Aksjomat}
\newtheorem{definition}[axiom]{Definicja}
\newtheorem{theorem}[axiom]{Twierdzenie}
\newtheorem{proposition}[axiom]{Stwierdzenie}
\newtheorem{lemma}[axiom]{Lemat}
\newtheorem{corollary}[axiom]{Wniosek}
\newtheorem{remark}[axiom]{Uwaga}
\newtheorem{hypothesis}[axiom]{Hipoteza}
\newtheorem{problem}[axiom]{Problem otwarty}
\newtheorem{result}[axiom]{Wynik numeryczny}
\newtheorem{chain}[axiom]{Krok}
\newtheorem{observation}[axiom]{Obserwacja}
\newtheorem{openproblem}[axiom]{Problem T-OP}

% --- Aliasy polskie (kompatybilność z plikami używającymi polskich nazw) ---
\newenvironment{twierdzenie}{\begin{theorem}}{\end{theorem}}
\newenvironment{lemat}{\begin{lemma}}{\end{lemma}}
\newenvironment{hipoteza}{\begin{hypothesis}}{\end{hypothesis}}
\newenvironment{propozycja}{\begin{proposition}}{\end{proposition}}
\newenvironment{definicja}{\begin{definition}}{\end{definition}}
\newenvironment{uwaga}{\begin{remark}}{\end{remark}}
\newenvironment{wniosek}{\begin{corollary}}{\end{corollary}}
\newenvironment{dowod}{\begin{proof}}{\end{proof}}
